% New Phase 10 fragment; no prior theory or numerical implementation is edited.
% Requires amsmath, amssymb, amsthm, hyperref and theorem/corollary environments.
\section{Finite-word memory forced by counterfactual admissions}
\label{sec:finite-word-admission-memory}

The continuous-input obstruction can be strengthened to a finite-alphabet
witness using exactly stored FP32 raw features and public binary labels.
The same raw features feed cosine curation and ridge learning; normalization
is confined to graph scoring.  The result is an information requirement for
prescribed ridge heads, not a total-byte or natural-language utility claim.

\subsection{Finite-format family and access contract}
Fix $b\ge1$, $m\ge2$, powers of four $k,D\ge1$ satisfying
$d=2+k+D\le2^{20}$, and a public rational $\lambda>0$.  Fix public signatures
\[
 u_i\in\{\pm1/\sqrt k\}^k,\qquad |u_i^Tu_j|\le1/6\quad(i\ne j).
\]
The theorem is conditional on this code property.  Orthogonal Hadamard rows
suffice when $m\le k$.  Normalized independent signs give existence if
$k\ge\lceil72\log(2m(m-1))\rceil$, by the pairwise bound
$2e^{-k/72}$ and a union bound; the chosen power of four must still satisfy
the total-dimension cap.  Arbitrary $m$ at fixed $k$ is not claimed.

Sort $N=1+b+2m$ fixed distinct public identifiers by the frozen hash priority,
including its ID tie break, independently of private data.  Assign roles in
that order to anchor $a$, common blockers $S=\{s_1,\ldots,s_b\}$, private
blockers $p_1,\ldots,p_m$, and candidates $w_1,\ldots,w_m$.  No assumption
about hash-collision absence is needed.  These are mathematical roles, not
synthetic empirical records or source-provenance claims.

Use orthogonal blocks $e_0,e_1,U,H$ of dimensions $1,1,k,D$.  Set
$\alpha=1/128$ and let $z_i\in\{\pm1/\sqrt D\}^D$ encode the only private
data: $mD$ independently varying signs.  Define the raw feature rows
\begin{equation}\label{eq:finite-raw-rows}
\begin{split}
 a&=e_0,&s_j&=\tfrac12e_0+\tfrac78e_1,\\
 p_i&=\tfrac12e_0+\tfrac78u_i,&
 w_i&=\tfrac12e_1+\tfrac78u_i+\alpha z_i.
\end{split}
\end{equation}
Candidate labels are one; all other labels are zero.  Identifiers, role map,
signatures, graph, labels, dimensions, regularization, and blocker payloads
are public.  Actual private signs/candidate rows are not.  Count every
sign-dependent model, cache, ticket, remote store, retained random seed,
and usable earlier output.  No retained-candidate reread or uncharged
data-dependent side channel is allowed.  Full forgotten payloads may be
supplied, but the distinguishing probes delete only fixed public blockers.
Public randomness is independent of the signs.

A $B$-bit state means at most $2^B$ distinguishable message values,
conditional on the public contract and public coins.  Fixed-capacity
$B$-bit storage satisfies this convention; any usable length, framing, or
allocation choice must be charged.  Strings of length at most $B$ have
$2^{B+1}-1$ possibilities if their length is free, and do not automatically
satisfy the convention.  The entropy proof uses $H(M\mid R)\le B$.

For selected $T\ne\varnothing$, the learner is the exact rational minimizer
of average-loss ridge on the \emph{stored raw} rows:
\[
 L_T(\theta)=\frac1{2|T|}\sum_{v\in T}(x_v^T\theta-y_v)^2
                          +\frac\lambda2\|\theta\|_2^2.
\]
For $T=\varnothing$ return zero.  This is not a bitwise FP32/FP64 solver
output contract.  The learner does not renormalize the stored rows.  All
claims are for record deletion; singleton-source deletion is formally
equivalent only when each record is its own owner, and does not establish
the empirical genuine-source arm.

\subsection{Uniform transfer to the frozen cosine graph}
The row norms are independent of the private signs:
\[
 \|a\|^2=1,\quad\|s_j\|^2=\|p_i\|^2=65/64,\quad
 \|w_i\|^2=16641/16384=(129/128)^2.
\]
All nonzero entries of \eqref{eq:finite-raw-rows} are exactly representable
normal FP32 values throughout the admitted dimension range.  The stored
threshold is the binary64 value
\[
 \tau=\operatorname{float64}(0.4)=3602879701896397/2^{53}
                                =2/5+1/(5\,2^{53}).
\]
The smallest required ideal edge cosine is
$448/(129\sqrt{65})>.43$, between a common blocker and a candidate.
The largest required nonedge cosine is at most
$18563/49923<.372$, between distinct candidates.  The latter includes the
worst private correlation $z_i^Tz_j\le1$; hence both bounds are uniform
over all $2^{mD}$ payloads.

The companion scorer proof analyzes the frozen ordered binary64
normalization and coordinate dot product.  Norm squares and partial sums
are exact in this family: each is a multiple of $2^{-34}$ of magnitude
below two.  With $u=2^{-53}$, its absolute score-error bound is
\[
 E_d=2\rho+\rho^2+\gamma(1+\rho)^2<2^{-32},\qquad
 \rho=\frac{2u}{1-u},\qquad
 \gamma=\frac{(d+1)u}{1-(d+1)u}.
\]
The numerical contract uses ordinary binary32/binary64, round-to-nearest
ties-to-even, correctly rounded square root/division, and the frozen
operation order without arithmetic-changing reassociation.  Nonzero
intermediates remain normal.  The frozen scorer does not clip.  This
arithmetic proof is not a general hardware certification or native-runtime
measurement.

These margins imply the same strict graph for every payload, with
\[
                          B(w_i)=S\cup\{p_i\}.
\]
Curator normalization leaves the raw learner values unchanged.  Using
normalized learner vectors would define a different target.

\subsection{Exact finite-alphabet state}
Write $t=16641/16384+2\lambda$ and $n=mD$.  Initially only the zero-label
anchor is selected.  At horizon $b$ no candidate can lose all $b+1$
blockers, so every target head remains zero.  At horizon $K=b+1$, the probe
$F_i=S\cup\{p_i\}$ selects exactly $\{a,w_i\}$; conversely every request
of size at most $K$ selecting a candidate equals its $F_i$.
Since $a\perp w_i$, its exact ridge head is
\begin{equation}\label{eq:finite-probe-head}
 \theta_i=\frac{w_i}{t}
   =\underbrace{\frac{\tfrac12e_1+\tfrac78u_i}{t}}_{\theta_i^{\rm pub}}
                          +\frac\alpha t z_i.
\end{equation}
The signs of its $H$ coordinates recover all $D$ private bits for candidate
$i$.

\begin{theorem}[Exact finite-word transition]
On this fixed public family the minimum deterministic exact model-query
state is zero private bits at horizon $b$ and $mD$ private bits at horizon
$b+1$.  The lower bound also applies to a sequential service's initial
state and to randomized services that are exactly correct almost surely
for every fixed input and query.
\end{theorem}
\begin{proof}
Two distinct sign assignments differ on some candidate word, and its fixed
public probe has different heads by \eqref{eq:finite-probe-head}.  Exact
state must distinguish all $2^{mD}$ assignments.  Storing the signs suffices:
reconstruct the indicated head on $F_i$ and return zero on other queries.
The low horizon always returns zero.  The randomized zero-error lower
bound also follows from the zero-distortion information bound below.
Each $F_i$ must be supportable as a sequential service's first request.
\end{proof}

This is private information, not total bytes: identifiers, public code,
signatures, graph, ledger, padding, and workspace still cost storage.
If $D\ge k$, for example $D=k$, then $D\ge(d-2)/2$; for fixed $b$ the
bound is proportional to $Nd$ on this admitted subfamily.  If most
dimensions are public, the exact statement remains $mD$, not $md$.
Fixed 32-bit private words consequently require at least
$\lceil mD/32\rceil$ words before additional metadata and workspace.

Keep the original horizon $K=b+1$.  For cumulative deletion $F$, let
\[
 I(F)=\{i:w_i\notin F,\ |B(w_i)\setminus F|\le K-|F|\}
     =\{i:F\subseteq S\cup\{p_i\}\},\qquad e(F)=|I(F)|.
\]
Empty continuation is allowed, so a currently selected candidate still
matters when the remaining deletion budget is zero.

\begin{theorem}[Conditional retained-state minimum]
Fix an admissible cumulative deleted set $F$ and demand correctness for
every private assignment and every allowed continuation.  Conditional on
this fixed public construction and ledger, the exact minimum is $e(F)D$
private bits, counting all usable sign-dependent auxiliary information.
\end{theorem}
\begin{proof}
Every $i\in I(F)$ remains isolatable by completing the deletion with
$B(w_i)\setminus F$.  Independently varying these words produces
$2^{e(F)D}$ distinguishable future-head families.  If $e(F)>0$, all records
in $F$ are fixed public blockers; their forgotten payloads reveal none of
these words.  For sufficiency store precisely the words indexed by $I(F)$.
Every current or future nonzero head is one of their probe heads.  Under
continuation $F'\supseteq F$, $I(F')\subseteq I(F)$, so discarded words
never become relevant within the fixed original horizon.
\end{proof}

The separate canonical codec provides concrete atomic updates, bit packing,
membership/horizon checks, and public-ledger accounting.  The lossy
query-summary upper bound below does not provide that sequential contract.
The theorem is a worst-case fixed-$F$ statement, not a claim that a
particular adaptive request policy reveals no information through its
realized requests.

\subsection{Expected squared parameter error}
For $0\le v\le1/4$ define
\[
 q(v)=\frac{1-\sqrt{1-4v}}2,\qquad
 R_{\rm sq}(v)=1-h_2(q(v)),
\]
and set $R_{\rm sq}(v)=0$ for $v\ge1/4$.

\begin{theorem}[Finite-word parameter information curve]
Let a state contain at most $B$ private bits, with public independent
randomness and fresh randomized repair allowed.  If every fixed input and
every fixed hard probe satisfies
$\mathbb E\|\widehat\theta_i-\theta_i\|_2^2\le\varepsilon$, then
\begin{equation}\label{eq:finite-rate}
 B\ge nR_{\rm sq}(v),\qquad
 n=mD,\qquad v=\frac{t^2\varepsilon}{4\alpha^2}.
\end{equation}
For this lower bound, average correctness under uniform independent signs
and a uniform index among the $m$ hard probes already suffices.
\end{theorem}
\begin{proof}
Write the private signs as $2Y_{ij}-1$ for uniform independent bits and
rescale the returned private coordinates:
\[
 \widehat Y_{ij}=\tfrac12\left(1+
                       \frac{t\sqrt D}\alpha\widehat\theta_i[H_j]\right).
\]
Clipping to $[0,1]$ cannot increase squared error, and before clipping
\[
 \frac1D\sum_j(\widehat Y_{ij}-Y_{ij})^2
 \le\frac{t^2}{4\alpha^2}\|\widehat\theta_i-\theta_i\|_2^2.
\]
For state $M$ and public coins $R$, posterior means
$p_l=\Pr(Y_l=1\mid M,R)$ minimize conditional squared error; their average
variance is at most $v$.  The function
$g(s)=h_2((1-\sqrt{1-4s})/2)$ is increasing and concave on $[0,1/4]$.
One direct check substitutes $s=q(1-q)$: the derivative sign reduces to
$-u+1/u+2\ln u\le0$ for $u\ge1$, since
$(u-1/u-2\ln u)'=(u-1)^2/u^2\ge0$.
Therefore conditional subadditivity and Jensen give
\[
 H(Y\mid M,R)\le\sum_l\mathbb E h_2(p_l)\le ng(v),\qquad
 B\ge I(Y;M\mid R)\ge n[1-g(v)].
\]
For $v\ge1/4$ the zero bound is immediate.
\end{proof}

The zero-information squared-error threshold is exactly
$\varepsilon_0=\alpha^2/t^2$: return $\theta_i^{\rm pub}$ on each hard
probe, giving this error for every payload.  Any fixed smaller error forces
a positive rate in \eqref{eq:finite-rate}.  Averaging over \emph{all} valid
deletion sets would be a weaker, unsuitable condition because zero-target
queries can dilute the loss.

\paragraph{Finite-length query-summary upper bound.}
Fix $0<v<1/4$, set $q=q(v)$, $r=\lfloor qn\rfloor$, $q_n=r/n$, and
\[
 V(n,r)=\sum_{j=0}^{r}\binom nj,\qquad
 K_{n,r}=\min\left\{2^n,
 \left\lceil\frac{2^n}{V(n,r)}(n\ln2+1)\right\rceil\right\}.
\]
A Hamming covering code of radius $r$ and size at most $K_{n,r}$ exists:
independent uniform codewords give expected uncovered population below
one at the displayed count; all $2^n$ words supply the alternative.
Fix a code and nearest-codeword rule publicly.  Draw an independent public
uniform mask and permutation, encode the masked/permuted word by its
covering-codeword index, and reproduce transformed coordinates as
$\widehat Y_\ell=q_n+(1-2q_n)c_\ell$, where $c_\ell$ is the codeword bit.
Undo the mask and permutation.
The coordinate-average squared error is at most
\[
 q_n^2+(1-2q_n)d_H/n\le q_n(1-q_n)\le v.
\]
The mask makes the transformed input independent of the permutation; thus
every fixed original coordinate has expected error at most $v$.  Reproduce
the public head component exactly and the hidden block from these soft
bits.  For every fixed input and hard probe this gives
\[
 \mathbb E\|\widehat\theta_i-\theta_i\|_2^2
 \le4\alpha^2v/t^2,\qquad B\le\lceil\log_2K_{n,r}\rceil.
\]
Non-hard queries return zero.  The endpoints use exact signs at $v=0$ and
the public-component predictor at $v=1/4$.

The type bound on $\binom nr$ gives the familiar
$nR_{\rm sq}(v)+O_v(\log(n+1))$ finite-length expression for fixed interior
$v$; the displayed finite formula is authoritative on the bounded format
domain.  The codebook and public randomness may be enormously expensive.
This is not an efficient, canonical, or sequentially erasable lossy codec,
and its guarantee is not for a query chosen after observing realized
state or coins.

\subsection{High-probability and certified-release bridges}
Fix $0\le\delta<1/2$.  If every fixed hard probe satisfies
$\|\widehat\theta_i-\theta_i\|_2\le\eta$
with probability at least $1-\delta$, a wrong coordinate sign costs at
least $\alpha^2/(t^2D)$ squared error.  The average bit error of sign
thresholding, or its better Bayes rule, is at most
\[
 p=\min\left\{\tfrac12,\ \delta+(1-\delta)
                              \min\{1,t^2\eta^2/\alpha^2\}\right\}.
\]
Binary Fano yields $B\ge n[1-h_2(p)]$.

More strongly set $a_H=\alpha/(t\sqrt D)$.  A successful certified radius
strictly below $a_H$ identifies the entire candidate word.  Block Fano gives
\[
 B\ge m\left[D-h_2(\delta)-\delta\log_2(2^D-1)\right]
 \ge m\left[(1-\delta)D-h_2(\delta)\right].
\]
At deterministic success this is the full $mD$-bit bound.  If
$0<\lambda\le1$, the dimension cap implies $t<4$ and $\sqrt D\le2^{10}$,
so $a_H>2^{-19}$; an actual certified parameter radius at most $2^{-20}$ is
sufficient.  This does not guarantee that a solver returns such a
certificate: moment error, solve error, and release rounding must all be
covered.  Distinct exact heads can otherwise collapse under low-precision
release; numerical zero and signed-zero bit patterns are distinct contracts.

\paragraph{Interpretation and attribution.}
The added qualification is finite-format realizability under the actual
frozen scorer and prescribed raw-feature ridge target, with a conditional
exact-state witness.  Distinguishability, bounded-deletion storage, and
rate-distortion tools are prior ideas.  The squared-Bernoulli derivation and
covering symmetrization specialize the earlier repository argument to
private feature bits; they are not a new general coding theorem.  See
Shannon, \emph{Coding Theorems for a Discrete Source With a Fidelity
Criterion} (1959), original-paper scan
\url{https://gwern.net/doc/cs/algorithm/information/1959-shannon.pdf},
and Scarlett and Cevher, arXiv:1901.00555v3,
\url{https://arxiv.org/pdf/1901.00555}.
Shannon's displayed binary Hamming curve is not the squared-soft-bit
formula derived here.  Tadakamalla, Surisetty and Asad (2026),
\emph{Predictive State Is Not Explanatory State}, also study query-relative
predictive/counterfactual memory; their official source release at commit
\texttt{4143f408843394e21a0622c9eb5d5de10a687a8c} implements recovery from
coordinate-flip threshold answers.  We do not claim that general
separation or partition-based sufficiency principle.  Its inaccessible
full paper leaves its approximate-success semantics unresolved in the
present literature pass.

The metric is parameter fidelity.  A public test law uniformly sampling
an $H$ basis vector has prediction error $\|P_H e\|_2^2/D$; an isotropic
unit-feature law over all coordinates introduces $1/d$.  Therefore the
constant parameter obstruction does not establish constant NLP task loss
or population excess risk.  The mathematical rows are not a demonstrated
semantic-encoder image.  No total-byte, native-runtime, or privacy-erasure
optimality follows.  The companion \texttt{SCORER\_TRANSFER.md} and retained
checker provide the arithmetic proof and bounded examples; separate codec
and independent proof checks retain their stated scopes.
